\chapter{CONCLUSIONES Y RECOMENDACIONES}

El presente capítulo consolida el cierre de la investigación según la fase de Conclusión del marco de \textit{Design Science Research} (DSR). Se exponen las conclusiones derivadas del diseño, modelado e implementación del protocolo de votación electrónica post-cuántico, y se proponen las recomendaciones pertinentes para futuras investigaciones en esta línea de estudio.

\section{Conclusiones}
% TODO: Escribir las conclusiones detalladas ligadas a cada objetivo específico 
% y al cumplimiento general de la hipótesis.
A continuación, se listan las principales conclusiones obtenidas de este trabajo de investigación:
\begin{itemize}
	\item \textbf{Conclusión sobre el diseño y fundamentación:} Se logró estructurar y fundamentar el protocolo criptográfico basado en los estándares FIPS 203 y FIPS 204.
	\item \textbf{Conclusión sobre el análisis formal:} El modelado formal simbólico demostró la robustez de las propiedades de integridad, confidencialidad y autenticidad del protocolo frente a un atacante con capacidades cuánticas.
	\item \textbf{Conclusión sobre el análisis empírico y rendimiento:} La implementación y simulación permitieron validar la viabilidad práctica del protocolo al comparar su latencia y comprobar que se mantiene dentro de los márgenes admisibles en relación a esquemas tradicionales.
\end{itemize}

\section{Recomendaciones}
% TODO: Escribir recomendaciones técnicas y metodológicas para futuros desarrolladores
% o investigadores interesados en expandir el alcance de este protocolo.
Con base en el desarrollo y las limitaciones del presente trabajo, se plantean las siguientes recomendaciones:
\begin{itemize}
	\item Ampliar el modelo formal para evaluar otras propiedades avanzadas del voto electrónico, como la verificabilidad universal o la resistencia a la coerción.
	\item Realizar pruebas de rendimiento sobre dispositivos móviles o de bajo consumo para evaluar su viabilidad en entornos heterogéneos.
	\item Estudiar la integración de esquemas híbridos que combinen criptografía clásica y post-cuántica en las fases de transición tecnológica de la seguridad informática.
\end{itemize}
